\section{Receiver-independent evaluation of weighted orbit folds}
\label{sec:lazy-weight-theory}

The existing native suffix-fold theorem~\cite{suffix} gives
$R\le R_0+2\tau$ canonical weight runs after $\tau$ truncations,
including signed vector weights. Static contraction cannot increase this
number. The mass bound and polynomial selector accounting are also already
established in the repository~\cite{weighted}. Classical weighted orbit computation was
proved polynomial by Agol--Hass--Thurston
\cite[Lemma~15 and Theorem~16]{aht}.
Our change concerns how much of that bounded state must be visited after
one particular fold. The preceding native implementation still scans the
entire run list, and performs vector arithmetic on every receiving run.
A lazy additive run tree removes both dependencies.

\subsection{Exact lazy actions on canonical run trees}

A run is a pair $(\ell,v)$, where $\ell$ is a positive integer length and
$v\in\mathbb Z^d$. A canonical run sequence has unequal values in consecutive
runs. Its lengths sum to the current universe size. Store it in a balanced
binary tree whose leaves are runs. Each subtree stores its leaf count,
height, total length, first value, last value, and a pending additive tag.
Heights are balanced by the number of leaves, independently of their
possibly enormous integer lengths.

\begin{lemma}[Uniform additions preserve internal run boundaries]
\label{lem:lazy-translation}
For any $u,v,z\in\mathbb Z^d$, one has $u=v$ if and only if
$u+z=v+z$. Consequently adding the same vector $z$ to every point of a
canonical subtree changes none of its internal run boundaries. It can be
represented by adding $z$ to the subtree tag and to its two endpoint
values. Canonicalization after reconnecting this subtree is required only
at the two external seams.
\end{lemma}
\begin{proof}
Subtract $z$ from the equality in either direction. The internal adjacencies
are unchanged, so every old unequal pair remains unequal. Only an adjacency
between an updated value and an unupdated value can change its equality
status, and these are precisely the external seams.
\end{proof}

The tag denotes a common increment not yet propagated to descendants.
Before a split or rotation changes a tagged node's descendants, the tag is
pushed to the current children and cleared at that node. Otherwise the scope
of the increment would change. This push rule is part of the invariant,
not an optional implementation convention.

\begin{proposition}[Run-tree primitives]
\label{prop:run-tree}
A deterministic height-balanced run tree on $R$ canonical runs supports
splitting at an arbitrary integer position, concatenation with boundary
coalescing, and adding a constant vector on a contiguous position interval
using $O(\log(R+2))$ structural operations and $O(d\log(R+2))$ integer
additions or comparisons. These operations need not inspect the other leaves
of the affected interval. Reporting or deleting $s$ specified consecutive
runs costs $O(s+\log(R+2))$ structural operations, apart from output work.
\end{proposition}
\begin{proof}
Use a leaf-oriented AVL tree, with empty-tree height zero and leaf height
one. Joining trees of heights $a,b$ takes $O(|a-b|+1)$ structural operations,
and its output height $c$ lies in
$\{\max(a,b),\max(a,b)+1\}$. Indeed, join descends the taller spine until
the two heights differ by at most one. The height gap decreases at each
descent, and each returning frame requires at most one single or double
rotation. Induction gives the stated output-height property. The minimum
number of leaves in an AVL tree satisfies the Fibonacci height recurrence,
so every height in question is $O(\log(R+2))$.

Stored subtree lengths identify the child containing a split position.
Only one root-to-leaf path is followed; if the position is inside a leaf,
that leaf becomes two positive-length runs of equal value in the two
separate output trees. Push pending tags along the path and use AVL joins
to reassemble its off-path pieces. The sum of these join costs is
logarithmic, not the product of two logarithmic bounds. To see this
explicitly, first observe inductively that each output fragment has height
at most that of the subtree being split. At a right turn on the split path,
let $a$ be the left sibling's height and $b$ the height of the left fragment
returned from the right child. The original AVL inequality gives
$b\le a+1$. The newly joined left fragment has some height $c\ge b$.
If $a\ge b$, then $a-b\le c-b$; if $b>a$, then $b-a\le1$.
Consequently this join costs $O(c-b+1)$. As the split returns toward the
root, the height of the accumulated left fragment increases only at right
turns. These height increases telescope to at most the original root
height. Left turns give the symmetric telescoping bound for the right
fragment. There are at most one root-height's worth of returning frames,
so all joins and tag pushes together take $O(\log(R+2))$ structural work.

For concatenation, compare the last realized value of the left tree with
the first realized value of the right. If unequal, use a balanced join.
If equal, remove these two extreme leaves, replace them by one leaf with
their combined length, and perform balanced joins. No new interior equality
can appear: the old neighbors of each removed leaf were unequal to its
value. Access, removal, and joining follow logarithmic-height paths.

For a range addition, split at its two endpoints, apply the tag update of
Lemma~\ref{lem:lazy-translation} to the middle tree, and concatenate the
three pieces canonically. Reporting a detached subtree visits each of its
leaves once. Empty pieces are omitted. All endpoint arithmetic is on binary
integers and never iterates over a leaf's numerical length.
\end{proof}

\subsection{Translation folds from one anchor and their jump locations}

Consider a legal native truncation from $[0,N)$ to $[0,C)$. For a positive
carrier let $p>0$ be its period and $b=C-p\ge0$. The deleted suffix is
pushed by
\[
 \rho(x)=b+((x-b)\bmod p)
\]
onto $[b,C)$. Let $f$ be the old suffix weight, extended by zero off
$[C,N)$, and define the added residue sum
\[
 F(r)=\sum_{x\equiv r\pmod p}f(x),\qquad b\le r<C.
\]
For $b<r<C$, finite summation gives
\begin{equation}\label{eq:lazy-residue-jump}
 F(r)-F(r-1)=\sum_{x\equiv r\pmod p}\bigl(f(x)-f(x-1)\bigr).
\end{equation}
One anchor determines the remaining additive constant. If the suffix
intersects $S$ constant runs $[l_j,h_j)$ with values $v_j$, then
\begin{equation}\label{eq:lazy-residue-anchor}
 F(b)=\sum_{j=1}^{S}
 \left(\left\lfloor\frac{h_j-1-b}{p}\right\rfloor
       -\left\lfloor\frac{l_j-1-b}{p}\right\rfloor\right)v_j.
\end{equation}
Only the clipped suffix run endpoints are used. Every left endpoint
contributes $+v_j$ to its residue jump and every right endpoint contributes
$-v_j$. Combine coincident locations, omit the jump at $b$ when integrating
from the anchor, and sort the remaining locations. This produces $O(S+1)$
constant addition intervals. The operations on quotient and residue values
use binary division and multiplication; they do not enumerate a residue
class or repeat a shift a numerical number of times.

\begin{theorem}[Receiver-independent native fold]
\label{thm:lazy-fold}
Suppose the old weight function has $R$ canonical runs and the deleted suffix
meets $S$ runs. A native translation fold can be evaluated exactly in
\[
 O\bigl((S+1)\log(R+2)+S\log(S+1)\bigr)
\]
structural and endpoint-order operations and
$O(d(S+1)\log(R+2))$ vector-coordinate operations, with binary integer costs
charged separately. A reflection fold obeys the same bound without the
$S\log(S+1)$ sorting term. Neither bound contains the number of receiving
runs or of untouched prefix runs. The resulting tree denotes the same
canonical function as the established native suffix-only operation.
\end{theorem}
\begin{proof}
Split off and report the deleted suffix, costing $O(S+\log(R+2))$ by
Proposition~\ref{prop:run-tree}. For translation, compute
\eqref{eq:lazy-residue-anchor} and the $O(S)$ residue events in
\eqref{eq:lazy-residue-jump}. Sorting and integrating them gives the exact
piecewise-constant function $F$ on the receiving interval. Add each nonzero
constant piece with the run-tree range-add primitive. There are $O(S+1)$
pieces, and all are inside the retained prefix. Their interiors are disjoint,
so the sum of the range updates is exactly $F$. The old prefix remains
pointwise unchanged except for the prescribed additions. Exact seam
coalescing after every update leaves a canonical result.

For reflection $x\mapsto K-x$, the image of $[l_j,h_j)$ is
$[K-h_j+1,K-l_j+1)$. These intervals occur in reverse order and remain
disjoint, so they can be reported directly in increasing image order and
added as constant ranges. No sorting of image endpoints is needed.

During successive range additions there can be temporary boundaries that
cancel only after a neighboring interval is updated. At most two boundaries
per update are introduced. Thus the transient run count is $O(R+S+1)=O(R+1)$;
using its logarithm instead of $\log(R+2)$ changes only the constant factor.
This justifies the height bound throughout the operation, not merely at its
end. The proof makes no assumption that a nonzero vector has positive
coordinates.
\end{proof}

\subsection{Full-trace cost and integer sizes}

Let a complete forward trace contain $\tau$ truncations and $E$ events in
all. For truncation $t$, let $S_t$ be its deleted suffix run count. For
contraction $j$, let $G_j$ be its number of static gaps and $D_j$ the number
of constant run pieces actually emitted. Put $R_* = R_0+2\tau$ using the
existing native run-growth theorem.

\begin{corollary}[Weighted replay ledger]
\label{cor:lazy-replay}
After canonical input preparation, the new weight evaluator requires
\[
 O\left(
 R_0+E+\sum_{t=1}^{\tau}
       \bigl((S_t+1)\log(R_*+2)+S_t\log(S_t+1)\bigr)
 +\sum_j\bigl(G_j\log(R_*+2)+D_j\bigr)
 \right)
\]
structural work, together with the corresponding exact vector arithmetic
and histogram costs. This is additional to validation and maintenance of
the unweighted pairing trace. It contains no cumulative receiver-run term
$\sum_t P_t$.
\end{corollary}
\begin{proof}
Sum Theorem~\ref{thm:lazy-fold}. For a contraction, split out each static gap,
report its runs to the histogram, remove it, and concatenate the surviving
pieces. Using position lengths rather than absolute labels automatically
closes the removed gap. Each gap costs logarithmic split/join work; reporting
costs exactly the deleted run pieces. The input tree can be built bottom-up
from its already canonical list in $O(R_0)$ work. Pairing deletion, trimming,
merger, and transmission do not change the weight function. No work for
those events is charged to the weight tree beyond the $O(E)$ event-dispatch
term. Pairing-row maintenance is charged separately.
\end{proof}

\Needspace{12\baselineskip}
For additive input intervals $[a_i,b_i)$ with vector values $v_i$, define
\[
 A=\sum_i(b_i-a_i)\lVert v_i\rVert_\infty.
\]
\begin{lemma}[Coefficient bound for delayed transfers]
\label{lem:lazy-bits}
During replay of legal native folds and contractions, each stored base
value, pending lazy tag, endpoint-value cache, and sum of inherited tags
has every coordinate bounded by $A$ in absolute value. The auxiliary
derivative and anchor arithmetic has magnitude $O(A)$. Hence replay uses
$O(1+\log(A+1))$ coefficient bits and $O(\log(N+1))$ endpoint bits, where
$N$ is the initial universe size. Input validation additionally reads all
supplied coefficient bits, including values assigned to empty intervals.
\end{lemma}
\begin{proof}
For the proof only, introduce one atom $(i,x)$ for every input-interval
occurrence $i$ and every integer $x\in[a_i,b_i)$, carrying vector $v_i$.
The sum of the atoms' infinity norms is $A$. No implementation step
constructs these atoms. The standard orbit-transport invariant partitions
them between active points and finalized orbits. A fold moves the atoms
at each deleted point to its retained image, and a contraction finalizes
the atoms at its deleted points. Thus the active allocations are disjoint.

Refine this mathematical invariant pointwise through the lazy tree.
Partition an active point's allocated atoms between its leaf's stored
base value and the pending tags along its root-to-leaf path. Each part's
vector sum equals the corresponding stored numerical value or tag.
The atom subsets attached to one constant-valued leaf or tag may differ
from point to point; their vector sums agree. Initially the complete
allocation is in the leaf base. A uniform addition transfers disjoint
atoms from the detached suffix into each receiving point; place them in
the updated tag, or in the base when the receiving subtree is a leaf.
Pushing a tag redistributes its pointwise subsets to the same children
before any child's scope changes. Splitting restricts these subsets to
the two intervals. A rotation changes only the bracketing after this
redistribution. Coalescing equal boundary leaves places each point's
current allocation in its new base. These operations preserve the
partition invariant. If a tag or addition has numerical value zero,
its zero-sum atom subset may be rebased pointwise without changing any
stored number; cancellation therefore causes no exception to the
invariant. This is auxiliary proof bookkeeping, not extra runtime work.

Every stored base, tag, inherited-tag sum, and realized endpoint cache is
therefore the sum of some subset of the input atoms, and its coordinate
magnitude is at most $A$. This also proves
$\sum_j(h_j-l_j)\lVert v_j\rVert_\infty\le A$ for the current runs.
An anchor's run multiplicity is at most its run length, so each partial
anchor sum has magnitude at most $A$. For the suffix extended by zero,
the sum of the norms of its endpoint jumps is at most
$2\sum_j\lVert v_j\rVert_\infty\le2A$. Every partial collision sum of
mapped jumps is therefore bounded by $2A$, and combining such sums with
an anchor is bounded by $3A$. All scalar multipliers in the anchor are
counts of integer points; all divisions use a strictly positive integer
period. Their operands have $O(\log(N+1))$ bits. The result follows.
\end{proof}

The fine structural bound is deterministic for the AVL transport and its
comparison sort of residue events. The delivered Python source
canonicalizer and histogram accumulator also use dictionaries. Their
map-operation counts are separate from the tree ledger: a constant-time
or average-case hash-table model must not be silently promoted to a
worst-case comparison bound. Sorting an endpoint-event list and using an
ordered histogram map or a sorted emission ledger would give deterministic
comparison bounds with the corresponding sorting factors. The delivered
code retains its established source and histogram routines.

The star-family separation below is a statement about the actual specified
flat and lazy evaluation methods. It is not a lower bound for every weighted
orbit algorithm, and it does not bound the number of topology hierarchy
states.

